\chapter{Consistency Across Producers, Brokers, and Consumers}

This chapter follows one message through three independent boundaries: producer confirmation and retry, broker replication and visibility, and consumer progress versus durable business effects. A successful result at one boundary is not an end-to-end transaction. The experiments use real TCP connections and per-broker logs; their observations distinguish physical log end (LEO), the high watermark (HW), consumer position, committed offset, and processor callbacks.

\lessonsection{29}{Cross-Layer Consistency Experiments}
\goals{Orchestrate reproducible failures across producer, broker, and consumer boundaries, then report what actually happened at each layer.}{Understand the difference between a successful append, a lost acknowledgment, HW visibility, a consumer's in-memory position, a committed group offset, and an external side effect.}
\background{The supplied \method{FaultProxy} forwards raw frames to a real broker and can either forward a matching request then drop its reply, or reject a matching request before forwarding it. \method{EffectRecorder} records every callback, including duplicates. \method{ConsistencyObservation} captures a named phase, error, replication snapshot, physical LEOs for all brokers, records fetched from the current leader below HW, optional consumer position and group committed offset, observed effects, and optional append receipt. \method{ConsistencyTrace} groups those observations for one partition. Build each scenario from the caller-owned, already-started \method{ClusterHarness}: an empty one-partition topic, leader 1, three replicas in ISR, minISR=2, HW=0, three physical LEOs at 0, no committed group offset, and unstamped input data. Fault scenarios also receive a started, unarmed proxy whose upstream is broker 1. Use the injected fixed clock and the supplied replication helper rather than sleeps or expiry races.}
Its API is \method{FaultProxy(Endpoint)}, \method{upstream()}, \method{start()}, \method{endpoint()}, \method{armed()}, \method{dropNextReply(short)}, \method{rejectNextRequest(short)}, \method{forwarded(short)}, \method{dropped(short)}, \method{rejected(short)}, and \method{close()}. Only one API fault can be armed at once. A drop fault forwards the request, reads the real broker reply, then closes the downstream connection without replying; a reject fault does not forward and returns \method{REQUEST_TIMEOUT} with \pathcode{ErrorBody("injected pre-forward failure")}. \method{LabSupport} supplies \method{replicateAndReport(cluster,tp)}, \method{observe(phase,error,cluster,tp,position,group,effects,receipt)}, and \method{awaitLeaderLeo(cluster,tp,expectedLeo)}. Observation fields, in order, are \pathcode{phase, error, replication, physicalLeo, position, committedOffset, visible, effects, receipt}; the visible records come from real leader FETCH below HW and physical LEOs come from each log, not authority reports.
\method{EffectRecorder} implements \method{RecordProcessor}; synchronized \method{process(tp,record)} appends every callback and synchronized \method{snapshot()} returns a stable copy without deduplicating. \method{ConsistencyObservation} is \pathcode{record ConsistencyObservation(String phase, ErrorCode error, ReplicationSnapshot replication, Map<Integer,Long> physicalLeo, OptionalLong position, OptionalLong committedOffset, List<LogRecord> visible, List<LogRecord> effects, Optional<AppendResult> receipt)}; \method{ConsistencyTrace} is \pathcode{record ConsistencyTrace(TopicPartition tp, List<ConsistencyObservation> observations)}.
\providededit{\method{FaultProxy}, \method{EffectRecorder}, \method{ConsistencyObservation}, \method{ConsistencyTrace}, and \method{LabSupport} provide fault injection, recording, and actual-state observation.}{\pathcode{exercises/src/main/java/io/simplekafka/lab/ConsistencyExperiments.java}: \method{producerOutcomeUnknown(ClusterHarness,TopicPartition,RecordData,FaultProxy)}, \method{leaderAckLoss(ClusterHarness,TopicPartition,RecordData)}, \method{allAckTimeout(ClusterHarness,TopicPartition,RecordData)}, \method{consumerReplay(ClusterHarness,TopicPartition,String,RecordData,EffectRecorder,FaultProxy)}, and \method{staleGroupSideEffect(ClusterHarness,TopicPartition,String,RecordData,EffectRecorder)}.}
Nonmatching APIs pass through unchanged; the proxy counters increase only for the corresponding real forwards, drops, or pre-forward rejections. An invalid API or an attempt to arm a second fault reports \method{INVALID_REQUEST}.
The five editable methods are public static and return \method{ConsistencyTrace}; \method{consumerReplay} also declares \pathcode{throws Exception}.
\lessoncontract{Implement all five methods using caller-supplied data. Require a started cluster with one empty partition, leader 1, ISR {1,2,3}, minISR=2, HW=0, all physical LEOs 0, and unstamped input; consumer scenarios also require an uncommitted group and an empty caller-owned recorder. The two fault-injection methods require a started, unarmed proxy whose upstream is broker 1. Invalid starting state reports \method{INVALID_REQUEST}. Close only scenario-owned clients; record only expected failures and propagate unrelated ones. \method{producerOutcomeUnknown} drops one successful LEADER reply and reports \pathcode{UNKNOWN} (LEO 1/0/0, HW 0, position 0, no visible record or receipt), then \pathcode{POISONED} when reuse of that \method{SimpleProducer} is rejected without another forwarded PRODUCE. An independent ordinary producer explicitly resends the same data and obtains [1,2); after replication report \pathcode{EXPLICIT_RESEND} (LEO 2/2/2, HW 2, position 2, visible offsets 0 and 1 contain equal input data, and no group committed offset is observed). \method{leaderAckLoss} obtains LEADER receipt [0,1) without replication and reports \pathcode{ACKED_NOT_VISIBLE} (LEO 1/0/0, HW 0, position 0); stop broker 1, elect broker 2 at epoch 1 and report \pathcode{ELECTED_WITHOUT_TAIL}; restart broker 1, reconcile, and report \pathcode{OLD_TAIL_REMOVED} (all LEOs 0, HW 0, no visible record). \method{allAckTimeout} uses ALL with timeout 0 and reports \pathcode{TIMEOUT_NOT_VISIBLE} (LEO 1/0/0, HW 0, position 0, no receipt); after real replication and reporting, the same consumer reads offset 0 and reports \pathcode{LATE_VISIBLE} (LEO 1/1/1, HW 1, position 1), without a resend or retroactive success ACK. \method{consumerReplay} writes and replicates one record to HW=1; the consumer reads offset 0 and the proxy rejects \pathcode{COMMIT_OFFSET} before forwarding. Report \pathcode{COMMIT_FAILED} with \method{REQUEST_TIMEOUT}, position 1, no committed offset, effects [0]. Close that consumer; a new connection resumes at 0, processes and commits again, then reports \pathcode{REPLAY_COMMITTED} with position and committed offset 1 and effects [0,0]. \method{staleGroupSideEffect} has member z join, fetch offset 0, and record an effect before member a joins and takes the assignment. z's old token gets \method{ILLEGAL_GENERATION} for both COMMIT and FETCH; report \pathcode{STALE_COMMIT_REJECTED} (position 1, no committed offset, effects [0]). The current owner fetches offset 0, records it, and commits nextOffset=1; report \pathcode{NEW_OWNER_COMMITTED} (committed offset 1, effects [0,0]). Fencing rejects stale requests but cannot undo an earlier callback.}
The \pathcode{UNKNOWN} and \pathcode{POISONED} observations carry \method{REQUEST_TIMEOUT}; the poisoned reuse is rejected locally and adds no proxy forward.
\lessonhints{Track the physical LEO of each broker separately from the authority's reported LEO and from HW.}{Write a timeline for each failure and label which side has already performed its action when the error occurs.}{Keep the recorder and proxy in the test caller: assert their counters and snapshots independently of the returned trace.}
\expected{Check every phase against the real cluster and caller-owned recorder/proxy. On response loss, expect one dropped PRODUCE reply and two forwarded PRODUCE requests: the poisoned attempt is not forwarded, and the explicit resend accounts for the second forward. In the election trace, \pathcode{ELECTED_WITHOUT_TAIL} is still LEO 1/0/0 at HW 0, then \pathcode{OLD_TAIL_REMOVED} is LEO 0/0/0. The ALL failure is \method{REQUEST_TIMEOUT}; its later visibility does not retroactively produce a receipt. For consumer replay, assert rejected(\pathcode{COMMIT_OFFSET})=1, forwarded(\pathcode{COMMIT_OFFSET})=0, effects [0,0], and a real \pathcode{FETCH_OFFSET} of 1 after the new consumer commits. For stale-group replay, the old token's two requests are fenced but effects remain [0,0].}
\pitfalls{A receipt is evidence of an ACK result, not proof that a LEADER-acknowledged tail will survive a later clean election. A timeout is not proof that no append occurred. A consumer position advances before a group commit, and group fencing protects broker operations rather than reversing a side effect. These experiments do not add retries to \method{SimpleProducer}, use sleeps to manufacture timeouts, or claim end-to-end exactly-once processing.}
\lessoncommand{29}
\lessonquestions{Which observations distinguish a LEADER acknowledgment from replication-committed visibility?}{Why can a rejected stale commit coexist with a side effect that has already happened?}

\lessonsection{30}{Idempotent Producer}
\goals{Persist producer batch identity with each record and make an explicit retry return the original offset range without appending duplicates.}{Understand batch sequence, producer epoch, content fingerprints, partial replication, and why retry behavior must be an explicit client choice.}
\background{The existing three-argument \method{RecordData} constructor remains the normal, non-idempotent record path. Its new nullable \method{producerStamp} component is metadata assigned by the broker for an idempotent request. \method{ProducerStamp} contains \method{producerId}, caller-managed \method{producerEpoch}, \method{firstSequence}, \method{batchSize}, \method{batchIndex}, and a 32-byte \method{batchHash}. The SHA-256 fingerprint is streamed over count and each record's timestamp, key length/key, and value length/value; it excludes offsets and leader epochs, and distinguishes null from empty keys. The fingerprint verifies that a retry with the same identity carries the same batch.}
A valid stamp has nonnegative producerId, producerEpoch, and firstSequence; positive batchSize; batchIndex in [0,batchSize); a non-overflowing \method{firstSequence+batchSize}; and exactly 32 hash bytes. Its fingerprint input is big-endian \pathcode{count:int32}, then per record \pathcode{timestamp:int64 | keyLength:int32(-1 for null) | key | valueLength:int32 | value}; offsets and leader epochs are excluded.

The supplied \method{ProducerStamp.ENCODED_OVERHEAD} is 64 bytes; that overhead counts against the 1 MiB record bound in both disk and \method{MessageCodec} wire budgets.
Each stamped record keeps the identity in its ordinary log entry, so FETCH, \pathcode{REPLICA_FETCH}, reconciliation, and disk recovery all see the same metadata. The stamped disk layout is \pathcode{length | crc32c | offset | timestamp | marker:int32=-2 | producerId:int64 | producerEpoch:int32 | firstSequence:int64 | batchSize:int32 | batchIndex:int32 | batchHash:32 bytes | keyLength:int32 | valueLength:int32 | key | value}. It adds 64 bytes, for a minimum total record size of 96 bytes; the 1 MiB record limit includes the stamp, and CRC32C covers offset through value. Ordinary unstamped bytes do not change. On the teaching wire, a stamped RecordData uses \pathcode{marker=-2 | producerId | producerEpoch | firstSequence | batchSize | batchIndex | batchHash[32] | keyLength | key | valueLength | value | timestamp}; ordinary wire records remain unchanged.

The broker derives deduplication state from the retained log prefix, not from a new authority truth. It validates sequences and complete batch hashes before writing, remembers completed batch ranges, and allows a replication fetch to append a matching partial batch suffix. A completed duplicate returns its original range; a retry that completes a partial tail appends only the missing suffix. A higher producer epoch starts at sequence zero and takes effect only with a durable record. A changed hash for the same batch identity conflicts; gaps, unrecorded old sequences, or invalid boundaries are out of order. An incomplete tail blocks unrelated appends until that original batch is completed or normal repair changes the prefix. If retained history starts above zero, producer state is expired rather than guessed.
\providededit{The stamp value, fingerprint helper, record/request serialization support, and replicated log infrastructure are supplied. Extend the existing storage and routing algorithms as well as implement the new client and replicated-produce path.}{\pathcode{exercises/src/main/java/io/simplekafka/client/IdempotentProducer.java}: \method{send(TopicPartition,List,Acks,long)}, \method{retry(TopicPartition)}; \pathcode{exercises/src/main/java/io/simplekafka/replication/IdempotentAppender.java}: \method{append(long,int,long,List)}, \method{recover()}; \pathcode{exercises/src/main/java/io/simplekafka/replication/ReplicatedPartition.java}: \method{produceIdempotent(List,Acks,int,long,long,int,long)}; \pathcode{exercises/src/main/java/io/simplekafka/broker/BrokerHandler.java}: \method{handle(IDEMPOTENT_PRODUCE)}; \pathcode{exercises/src/main/java/io/simplekafka/client/MetadataRouter.java}: \method{call}; extend \method{RecordCodec.encode/decode}, \method{SegmentLog.append/recover}, and \method{PartitionLog.append/deleteBefore/truncateTo}.}
The supplied partition-log support also exposes \method{tailProducerStamp()}, \method{prefixVersion()}, and \method{truncateToOffset(long)} so append validation and cache invalidation can follow the actual persisted tail and rewritten prefix.
The supplied \method{IdempotentProducer(MetadataRouter,long,int)} constructor and \method{close()}, and the \method{IdempotentAppender(PartitionLog)} constructor, provide lifecycle setup; implement only the listed algorithm methods.
\method{IdempotentAppender} construction does not automatically run \method{recover}; recovery is an explicit algorithm call.
The public algorithms return \method{ProduceReceipt} from \method{send(TopicPartition,List<RecordData>,Acks,long)} and \method{retry(TopicPartition)}, \method{AppendResult} from \method{IdempotentAppender.append(long,int,long,List<RecordData>)} and \method{ReplicatedPartition.produceIdempotent(List<RecordData>,Acks,int,long,long,int,long)}, and \method{void} from \method{recover()}.
The supplied \method{IdempotentProduceBackend} capability is implemented by \method{ReplicatedPartition}; \method{BrokerHandler} rejects API 11 with \method{INVALID_REQUEST} when a backend lacks it, including the single-broker backend.
\lessoncontract{\method{IdempotentProducer.send} fixes the per-partition sequence and a stable copy of the complete input batch before one RPC. Only a valid successful receipt advances the sequence and clears pending state. Any request/reply failure keeps the batch pending; another send to that partition is rejected, while \method{retry} explicitly refreshes metadata and makes exactly one request with the same identity, sequence, payload, ACK mode, and timeout. There is no background or automatic retry, and closing the producer discards only local pending state without closing its borrowed router. The caller supplies producer id and epoch; a rebuilt client must use a higher epoch and sequence zero. \method{IdempotentAppender} validates all inputs before mutation, returns the original range for an exact completed duplicate, appends only the missing suffix of an identical partial batch, and rejects a same-identity/different-hash batch with \method{DUPLICATE_SEQUENCE_CONFLICT}. Old producer epochs report \method{FENCED_PRODUCER_EPOCH}; gaps and invalid old sequences report \method{OUT_OF_ORDER_SEQUENCE}; retained-prefix loss reports \method{PRODUCER_STATE_EXPIRED}; unrelated writes while a tail batch is incomplete report \method{INCOMPLETE_BATCH}. After a durable mutation fails or leaves ambiguous state, the log is failed and cannot accept further operations until it is reopened. Prefix truncation or retention invalidates derived cache state; recovery scans and verifies the retained history. The wire API is id 11, \method{IDEMPOTENT_PRODUCE}, with request fields \method{tp, epoch, acks, timeoutMillis, producerId, producerEpoch, firstSequence, records}; success uses the existing \method{ProduceBody}. Reject empty, oversized, invalid, or already stamped client batches without side effects. This API is available only through the replicated backend, not the single-broker backend. Duplicate requests still pass ACK prechecks and the full ACK wait, including ALL's minISR, epoch, HW, and deadline rules.}
The log-derived state tracks each producer's highest durable epoch, next sequence, and completed batch identity/range; the partition has at most one incomplete tail batch. An unknown producer or a higher epoch starts at sequence zero; within the current epoch, a new batch must start at exactly the stored next sequence. Recovery reads in bounded batches, verifies batch continuity and hashes, and replays only the new suffix while the prefix version is unchanged. If the prefix changes or LEO rolls back, discard the cache and rebuild from offset zero; a recovery error invalidates the derived state rather than skipping damaged history.
Client details: a pending batch blocks only another \method{send} for the same partition; other partitions are independent. A successful receipt must cover exactly the batch's record count before sequence advancement. \method{retry} with no pending batch reports \method{INVALID_REQUEST}; refresh failure and invalid replies preserve pending state, and a retry uses the refreshed current leader epoch. A repeated batch still performs ACK precheck and the full ACK wait.
\lessonhints{Keep the three identities separate: leader epoch fences broker roles, producer epoch fences a producer session, and sequence identifies a batch within that producer session.}{Treat the log as the durable source of truth and the in-memory producer-state cache as a versioned replay result.}{Test one complete duplicate, one same-sequence changed-content conflict, one partial-batch continuation, and one explicit retry after a leader change; inspect both receipt ranges and stamped disk records.}
\expected{Send a multi-record batch and repeat the same identity and payload: both receipts name the same half-open range, while physical LEO and stamped bytes do not change on the duplicate. The following sequence appends at the next offset. After the first response is dropped, copy the records to HW, elect a clean new leader, and explicitly retry to obtain the original range with only one copy on every broker. Reopen a broker's log from disk and issue the same raw request: it still returns the original range. A repeated ALL request must still time out while its range is above HW and can succeed only after the range is replicated and acknowledged.}
\pitfalls{This teaching identity is caller-supplied; it does not implement Kafka PID allocation, a transaction coordinator, or client session WAL. A new producer id does not deduplicate a repeated business event, and a business event id is not a producer sequence. \method{SimpleProducer} remains non-idempotent and never retries; the new \method{IdempotentProducer} retries only when its caller invokes \method{retry}. Idempotent append does not strengthen ISR, HW, or election assumptions and does not create an end-to-end exactly-once guarantee.}
\lessoncommand{30}
\lessonquestions{Why must a duplicate ALL request still check the current HW and minISR?}{What information is lost if producer deduplication state cannot be reconstructed from the retained log prefix?}

\lessonsection{31}{Idempotent Consumer Effects}
\goals{Apply each stable business event to a durable local ledger once, even when its Kafka record is delivered more than once.}{Understand why delivery deduplication and business-event deduplication use different identities, and why external effects are not atomic with consumer offsets.}
\background{\method{BusinessEvent} contains a stable \method{eventId}, an \method{account}, and a signed \method{delta}. Both identifiers are non-empty strict UTF-8 strings of at most 255 bytes; delta may be positive, negative, or zero. Its versioned value encoding is \pathcode{version:uint8=1 | eventId:identifier | account:identifier | delta:int64}. The supplied \method{MessageCodec.encodeBusinessEvent/decodeBusinessEvent} uses the course identifier encoding and rejects malformed or trailing data. The store appends the encoded event id, account, and delta as one ledger record; its in-memory identity and balance maps are rebuilt by replay rather than persisted separately.}
The codec signatures are \pathcode{public static byte[] encodeBusinessEvent(BusinessEvent)} and \pathcode{public static BusinessEvent decodeBusinessEvent(byte[])}; they reject unknown versions, truncated input, invalid UTF-8 or identifiers, and trailing bytes.
For the consumer demo, the supplied processor decodes \method{record.data().value()} with \method{MessageCodec.decodeBusinessEvent} and calls \method{DurableEffectStore.apply}; \method{ProcessingLoop} keeps its existing process-then-commit order. The supplied \method{balance(account)} validates the account identifier by the same rule and returns zero when no balance exists.
\providededit{\method{BusinessEvent}, event serialization, and \method{DurableEffectStore(Path)}, \method{balance(String)}, and \method{close()} are supplied.}{\pathcode{exercises/src/main/java/io/simplekafka/client/DurableEffectStore.java}: implement \method{apply(BusinessEvent)} and \method{recover()}.}
The supplied store API is \pathcode{DurableEffectStore(Path)}, \pathcode{synchronized boolean apply(BusinessEvent)}, \pathcode{synchronized long balance(String)}, \pathcode{synchronized void recover()}, and \pathcode{synchronized void close()}; close is idempotent.
The supplied constructor opens the local log and rebuilds store state through \method{recover()} before use.
\lessoncontract{Use \pathcode{<directory>} as a local \method{PartitionLog}. One event is one \method{RecordData(null,encodedEvent,0)} and the record append is forced before publishing in-memory state. On the first valid event id, compute the new account balance with checked addition, append once, then publish the full event in \method{seen} and the new balance; return true. An identical event id, account, and delta returns false and changes neither log nor balance. Reusing an event id with different content or overflowing a balance reports \method{INVALID_REQUEST} without appending. \method{recover} rebuilds both maps in retained offset order and requires logStart=0, null keys, timestamp 0, and no producer stamp. It decodes each payload, validates event identities, and counts an identical repeated event only once; malformed complete records, conflicting ids, overflow, or a scan with no progress report \method{CORRUPT_RECORD}. Synchronize apply, balance, recover, and close on the same store instance. If an append fails with an uncertain durable result, publish no new in-memory balance and mark the store failed; later operations report \method{STORAGE_ERROR}, so the caller must close and reconstruct from disk.}
Recovery scans the retained ledger in bounded reads and applies records in offset order; an incomplete physical tail follows the underlying log's recovery rule, while complete malformed records are rejected.
A closed or failed store rejects \method{apply}, \method{balance}, and \method{recover} with \method{STORAGE_ERROR}; \method{close} remains idempotent.
\lessonhints{The event id must come from the business event and stay the same across Kafka offsets and process restarts.}{Validate the event-id conflict and checked balance before touching the log; publish memory only after the append succeeds.}{Replay the same event from a new consumer position after closing and reopening the store, then compare both its balance and its ledger records.}
\expected{Process event A (account +10), a second Kafka record carrying the same A, and event B (account -3). The first A and B each append once; the repeated A returns false, leaving balance 7 and ledger entries [A,B]. If the first commit fails, rebuild the consumer and replay offsets 0, 1, and 2; verify callback history [0,1,2,0,1,2], the same two ledger entries, balance 7, and committed offset 3. A stale generation can reject the commit but cannot reverse the durable effect.}
\pitfalls{The store makes one local ledger append atomic with its in-memory publication; it does not share a transaction with Kafka or the consumer offset store. A callback can run more than once, including for the same event at different offsets. Group fencing rejects an old token but cannot roll back an effect. This is application-level idempotent accounting for one local store, not a distributed transaction, global ledger, or end-to-end exactly-once claim.}
\lessoncommand{31}
\lessonquestions{Why is a stable event id a better deduplication key than a Kafka offset or a callback count?}{What remains uncoordinated when the ledger append succeeds but the consumer offset commit fails?}
